\documentclass[11pt]{article}
\usepackage[utf8]{inputenc}
\usepackage{amsmath,amssymb}
\usepackage{graphicx}
\usepackage{booktabs}
\usepackage{authblk}
\usepackage[margin=1in]{geometry}
\usepackage[hidelinks]{hyperref}
\usepackage{setspace}

\newcommand{\grad}{\nabla}
\newcommand{\half}{\tfrac12}

\title{RealQM and the Electron Gas:\\ what extends, what does not, and what a
zero-temperature pressure would require}

\author[1]{Claes Johnson}
\affil[1]{KTH Royal Institute of Technology, SE-100 44 Stockholm, Sweden \\
\texttt{claesjohnson@gmail.com}}
\date{\today \\[0.4em]
\small\textit{Mathematical formulation by the author; analysis developed} \\
\small\textit{with AI assistance (Claude, Anthropic).}}

\begin{document}
\maketitle

\begin{abstract}
\noindent RealQM models matter as non-overlapping unit charge densities in ordinary
three-dimensional space, one per electron, meeting at free boundaries and minimising the Coulomb
energy. It has been developed and tested on \emph{bound} matter --- atoms, molecules and nuclei ---
where an attractor is always present. This note asks what happens when the framework is carried to a
domain it was not built for: the electron gas, and with it plasmas, metals and dense matter. The
answer separates cleanly. The \textbf{form} of the theory extends exactly: partitioning a gas of
density $n$ into non-overlapping cells gives a kinetic energy density $Cn^{5/3}$ --- the
Thomas--Fermi form, obtained here from non-overlap rather than from Fermi statistics --- and
everything that follows from charge and inertia is reproduced without adjustment, the plasma
frequency $\omega_p^2=4\pi n$ exactly. What does \emph{not} extend is the coefficient: for a uniform
gas with the free-boundary interface used throughout RealQM, the minimising density is constant, so
$C=0$. The model then has no pressure at zero temperature, hence no screening length and no plasmon
dispersion, so the zero-temperature kinetic energy of matter comes entirely from density variation forced by
attractors and nothing from the partition itself. A perfectly uniform gas is an idealisation and no
real system is one; what the result establishes is that the \emph{mechanism} is absent, not that
sodium has no stiffness. Whether lattice-induced modulation supplies enough in a real metal is an
open computation, and the bulk modulus is the place to settle it. Thermal motion does not help:
it supplies a pressure $\sim nk_BT$ which vanishes as $T\to0$, while the missing one does not. The obvious repair --- a Robin interface
condition $\psi'+\beta\psi=0$ --- reproduces Thomas--Fermi exactly at the single dimensionless value
$\beta a=1.146$, but must then be applied to atoms as well, where the interface passes through the
nucleus and the density is greatest --- so atomic energies should be strongly sensitive to it. A
first numerical test of that expectation was not converged and is reported as inconclusive. We conclude that extending RealQM beyond bound matter requires a genuinely new
ingredient: a mechanism by which compressing a partition costs energy at zero temperature. That is
stated here as a well-posed open problem rather than a defect, since it lies outside the domain in
which the framework has so far been tested.
\end{abstract}

\setstretch{1.15}

\section{Why the electron gas is a new domain}

RealQM \cite{realqm,realqmbook} represents an $N$-electron system by $N$ non-negative densities
$\psi_i$ on non-overlapping domains $\Omega_i$, with energy
\begin{equation}\label{eq:E}
E[\{\psi_i\},\{\Omega_i\}]=\sum_i\int_{\Omega_i}\Bigl(\half|\grad\psi_i|^2
-\sum_a\frac{Z_a}{|x-x_a|}\psi_i^2\Bigr)dx
+\half\sum_{i\neq j}\iint\frac{\psi_i^2(x)\psi_j^2(y)}{|x-y|}\,dx\,dy,
\end{equation}
subject to $\int_{\Omega_i}\psi_i^2=1$, the interfaces free to move to where the energy says. There
is no wavefunction in $3N$ dimensions, no antisymmetry and no self-interaction; the role played by
the exclusion principle is played instead by the requirement that the domains do not overlap.

Every test of this construction so far --- ionisation energies across the periodic table, molecular
binding, nuclear binding by dual Coulomb confinement --- shares one feature: an \emph{attractor} is
present. A nucleus is always there to shape the density. The question this note asks is what happens
when it is not, which is the situation in an electron gas, a plasma, or a metal, and which is a
domain RealQM was neither built for nor tested on. It should not be surprising if new ingredients
are required, and the purpose here is to identify exactly which.

\section{The form extends: non-overlap gives the Thomas--Fermi scaling}

Partition a gas of density $n$ into non-overlapping cells, one per electron, of linear size
$a=n^{-1/3}$. The localisation cost of a cell scales as $\hbar^2/2m\,a^{-2}\propto n^{2/3}$, so the
kinetic energy density is
\begin{equation}\label{eq:tf}
t(n)=C\,n^{5/3},
\end{equation}
with $C$ fixed by the shape of the cell and the condition at its faces. This is precisely the
Thomas--Fermi form, and it is worth pausing on where it came from: not from Fermi statistics, but
from the geometry of a partition. The framework's central premise --- that exclusion is a statement
about regions rather than about exchange symmetry --- delivers the correct \emph{scaling} of the
degenerate electron gas immediately.

Everything downstream then follows from $C$. Minimising \eqref{eq:E} at fixed chemical potential
with a test charge $Q$ present, and linearising about the uniform state, gives
$\delta n=(9n_0^{1/3}/10C)\varphi$ and hence, with Poisson's equation, a Helmholtz problem
$\grad^2\varphi-k^2\varphi=-4\pi Q\delta(x)$ with
\begin{equation}\label{eq:k2}
k^2=\frac{18\pi\,n_0^{1/3}}{5C},
\end{equation}
so that $\varphi=Q\,e^{-kr}/r$: Yukawa screening, with $\lambda=1/k\propto n^{-1/6}$. Setting $C$ to
the Thomas--Fermi coefficient $C_{\rm TF}=\tfrac{3}{10}(3\pi^2)^{2/3}=2.8712$ reproduces
$k_{\rm TF}^2=(4/\pi)(3\pi^2n)^{1/3}$ \emph{exactly}, at every density (Table~\ref{tab:lambda}).

\begin{table}[h]\centering
\small
\begin{tabular}{@{}lrr@{}}
\toprule
$n$ (a.u.) & $\lambda$ from \eqref{eq:k2} with $C_{\rm TF}$ & $\lambda_{\rm TF}$\\
\midrule
$0.01$ & $1.086\,a_0$ & $1.086\,a_0$\\
$0.1$  & $0.740\,a_0$ & $0.740\,a_0$\\
$1.0$  & $0.504\,a_0$ & $0.504\,a_0$\\
\bottomrule
\end{tabular}
\caption{The form is exact. Everything rests on the single constant $C$.}
\label{tab:lambda}
\end{table}

\section{What extends without adjustment: the plasma frequency}

Displace the electron distribution rigidly by $\xi$ against the ion background. The resulting surface
charge produces a restoring field $4\pi n\xi$, so $\ddot\xi=-4\pi n\xi$ and
\begin{equation}
\omega_p^2=4\pi n,
\end{equation}
free of temperature and of any statistical assumption. \textbf{RealQM reproduces this exactly}, and
for a structural reason: under a rigid translation each $\psi_i$ is carried along unchanged, so
$\grad\psi_i$ is unchanged, the localisation term contributes nothing, and what remains is Coulomb
and inertia --- which is all the result depends on. No property of $C$ enters.

The same argument identifies what RealQM should do about \textbf{Landau damping}. The framework is
not a single fluid: $N$ electrons are $N$ domains, each carrying its own current and hence its own
drift velocity, which is structurally a multi-beam plasma. Multi-beam models reproduce Landau damping
as phase mixing among cold streams, and do so reversibly, in agreement with plasma-echo experiments.
Where the velocity spread is thermal, the framework should therefore give collisionless damping for
the right reason rather than by an inserted term.

\section{What does not extend: the coefficient is zero}

The step from \eqref{eq:tf} to a number requires the cell to have a nonzero localisation cost, and
with RealQM's own interface condition it does not.

The kinetic term in \eqref{eq:E} is integrated over each electron's \emph{own} domain, and no node is
imposed where domains meet; the calibration of atomic RealQM found the free (Neumann) plane decisive
against a Dirichlet node. Take then $\psi_i\equiv|\Omega_i|^{-1/2}$, constant on each cell. This
satisfies the normalisation, makes the total density $\sum\psi_i^2$ uniform, minimises the
electrostatic energy against a neutralising background, and has $\grad\psi_i=0$ identically. The
kinetic energy is exactly zero, and the discontinuity at the cell face costs nothing because the
integral runs over $\Omega_i$ alone.

The uniform state is therefore the global minimum, with
\begin{equation}
C=0
\end{equation}
at every density. By \eqref{eq:k2} the screening length is infinite. The compressibility correction
to the plasma frequency, $\omega^2=\omega_p^2+\tfrac35k^2v_F^2$, vanishes, so the predicted plasmon
is dispersionless. And the pressure of a cold electron gas is zero.

Two attempts to locate a missing term fail. Allowing the interfaces to move changes nothing: with
constant $\psi_i$ the kinetic energy remains zero and the electrostatics merely fixes the cells to
equal volume. Requiring each domain to be connected and to carry unit charge changes nothing either,
since any equal-volume connected tiling admits the same constant solution. The result is a property
of the functional, not an oversight in its evaluation.

\section{The interface condition cannot repair it}

Since the Neumann condition gives $C=0$ and a hard wall gives $C=3\pi^2/2=14.804$, the two bracket
the target, which suggests a Robin condition $\psi'+\beta\psi=0$ at the interface. For a cubic cell
of side $a$ the problem separates, with $\kappa=ka$ solving $\kappa\tan(\kappa/2)=\beta a$ and
$C=\tfrac32\kappa^2$:

\begin{table}[h]\centering
\small
\begin{tabular}{@{}lrrr@{}}
\toprule
interface condition & $\kappa$ & $C$ & $C/C_{\rm TF}$\\
\midrule
Neumann, $\beta a=0$ (present choice) & $0$ & $0$ & $0.00$\\
Robin, $\beta a=1.146$ & $1.3835$ & $2.8712$ & $\mathbf{1.00}$\\
Dirichlet, $\beta a\to\infty$ & $\pi$ & $14.804$ & $5.16$\\
\bottomrule
\end{tabular}
\caption{One dimensionless number reproduces Thomas--Fermi exactly --- and with it the screening
length, the plasmon dispersion and the zero-temperature pressure, since all three follow from $C$.}
\label{tab:robin}
\end{table}

This is not adopted here, but the reason must be stated more carefully than a first calculation
suggested. A Robin condition at $\beta a=1.146$ would have to be applied everywhere, including
between the electron domains of an atom, and there is a geometric reason to expect atoms to object:
in helium the two domains are half-spaces meeting at a midplane, and \emph{the nucleus lies in that
plane}, where the density is greatest. Any departure from Neumann suppresses density exactly where
binding is won, so atomic energies should be strongly sensitive to $\beta$ --- unlike a uniform gas,
where the density is constant and the interface condition is the only place physics could enter.

That expectation has \emph{not} been confirmed numerically here, and the attempt is reported because
its failure is instructive. Solving helium as two half-space domains on a uniform grid, with
imaginary-time relaxation and the interface condition varied, gave a shift of $+0.19$~Ha at $N=64$
and $-0.43$~Ha at $N=96$ --- reversing sign under refinement, with the total energy moving from
$-2.02$ to $-0.94$~Ha against an exact $-2.9037$. The cause is that the relaxation step scales as
$h^2$ while the iteration count was held fixed, so finer meshes received less relaxation and were
further from the minimum. Neither the magnitude nor the sign of the sensitivity is therefore
established, and no conclusion is drawn from it.

What remains is the structural asymmetry, which does not depend on that calculation.
In a uniform gas the density is constant, so
the interface condition is the only place physics could enter --- and Neumann makes it contribute
nothing. In an atom the interface passes through the density maximum, so the condition is decisive
--- and only Neumann leaves the atom alone. \emph{The interface is inconsequential exactly where the
gas needs it to matter, and decisive exactly where the atom needs it not to.} No single $\beta$
serves both, and a value that switched between regimes could not be applied at all, since partially
ionised plasmas and metals contain bound and free electrons in the same system with no surface
between them.

\section{The role of temperature, and why it does not close the gap}

At finite temperature the domains move and fluctuate in shape, the density within a cell is no longer
constant, and $\grad\psi_i\neq0$: there is a kinetic energy, and with it a pressure. RealQM should
reproduce it, since it follows from charge, inertia and thermal motion, all of which the framework
carries. But this pressure scales as $nk_BT$ and vanishes as $T\to0$.

The missing pressure does not. Splitting the two contributions:

\begin{table}[h]\centering
\small
\begin{tabular}{@{}lcc@{}}
\toprule
contribution & survives $T\to0$? & RealQM\\
\midrule
thermal, ${\sim}nk_BT$ & no & should be reproduced\\
degenerate, ${\sim}nE_F$ & \textbf{yes} & \textbf{zero}\\
\bottomrule
\end{tabular}
\end{table}

\noindent and the second is not an exotic regime. A metal at room temperature has
$T_F\sim10^4$--$10^5$~K, so $T/T_F\approx10^{-2}$: sodium and aluminium are held up almost entirely
by the pressure that survives at zero temperature. The conflict is therefore blunt and requires no
appeal to any formalism: \emph{the model predicts that a cold dense electron gas exerts no pressure,
and sodium exists}. The same statement applies to white dwarf support and to the interiors of dense
stars.

It is worth being clear about what is and is not being claimed. Antisymmetry of a many-body
wavefunction is a formalism and is not itself observed; the framework is entitled to account for the
phenomena by other means, and that is its programme. What is observed --- the mass--radius relation
of white dwarfs, the linear electronic specific heat, de Haas--van Alphen oscillations, the bulk
moduli of simple metals, positive plasmon dispersion --- is that a large pressure exists at
temperatures far below $T_F$. A theory of matter must produce it somehow.

\section{Electric conduction: the same boundary, seen a second time}

Conduction in a metal is the interaction of an electron gas with a lattice of positive ions, and it is
worth asking separately, because the lattice changes one thing in the framework's favour and leaves
two others against it.

In favour: a \emph{gitter} of ions supplies attractors throughout the material, so the density is
shaped everywhere and the mechanism of \S5 does not apply --- $C=0$ was a statement about a
\emph{uniform} gas. Whatever else is true, a metal is not the case where the localisation term
vanishes identically.

Against, first: metallic conduction is the paradigm of \emph{delocalisation}. The carriers are
extended states spanning the crystal, whereas RealQM's electrons are localised by construction, each
confined to its own domain, with non-overlap requiring $N$ carriers to occupy $N$ disjoint regions.
Interpenetrating extended domains are not excluded in principle, but the natural minimiser of
\eqref{eq:E} is a set of compact cells. A theory of that shape describes a Mott insulator more
naturally than a conductor.

Against, second, and more sharply: the distinction between a metal and an insulator is in standard
theory a matter of \emph{band filling} --- two electrons per band per cell, and whether the count
leaves a band partly filled. That counting is a consequence of the exclusion principle, and RealQM
has no bands and no such count. Nor is there a Fermi surface, so the carrier density in
$\sigma=ne^2\tau/m$ has no evident referent: standard theory counts only the electrons within
${\sim}k_BT$ of $E_F$, and here there is no $E_F$ to be near.

What the framework does look suited to is the opposite transport regime. Localised charges on sites,
moving by thermally activated jumps between domains, is \emph{hopping} conduction --- the mechanism
of semiconductors, disordered systems and small-polaron transport. That requires no Fermi surface,
sits naturally with non-overlapping domains, and would emerge from the same thermal motion that
supplies the pressure of \S6.

This yields a prediction crude enough to be decisive, and it concerns only a sign. Metallic
conductivity \emph{falls} with temperature, as thermal disorder scatters extended carriers; hopping
conductivity \emph{rises} with temperature, as jumps are activated. If RealQM transports charge by
hopping and not by band motion, it should predict $d\sigma/dT>0$ for every material --- that is, it
should make copper behave like a semiconductor. Whether that is what the framework actually gives has
not been computed here, and it is the natural first test of RealQM in a solid.

The pattern is the same one as before: bound and localised, well described; free and delocalised,
not yet.

\section{What an extension would require}

The open problem can be stated in the framework's own terms, without reference to exchange symmetry:

\begin{quote}
\emph{What geometric property of a partition should cost energy when its cells shrink?}
\end{quote}

Non-overlap alone does not, because a constant density in a cell has no gradient however small the
cell. Yet compressing a metal demonstrably costs energy. Something in the description of a partition
must resist compression at zero temperature, and identifying it is what carrying RealQM beyond bound
matter requires.

Two remarks may help locate it. First, the failure is specific to the \emph{uniform} case: wherever
an attractor shapes the density, the profile itself supplies the gradients, which is why the atomic
and molecular results are unaffected and why the deficiency was not visible before. Second, the
correct coefficient lies between the two principled interface conditions and is produced by neither
($0$ and $5.16$ against a required $1.00$), which is the signature of a missing mechanism rather than
of a mis-set parameter.

Until such a mechanism is found, the honest statement of scope is that RealQM is a theory of
\textbf{bound} matter --- atoms, molecules and nuclei, where an attractor is present --- and that its
extension to free-electron matter awaits an ingredient it does not yet contain. Nothing in the
established results depends on the extension, and stating the boundary costs the framework nothing it
had.

\section*{A note on what was tested numerically}

The screening relation \eqref{eq:k2}, the vanishing of $C$, the plasma frequency and the Robin table
are analytic. Table~\ref{tab:he} is a direct grid solution of helium as two half-space domains with
the interface condition varied, and its absolute energies ($-2.02$~Ha against the experimental
$-2.9037$) reflect a coarse mesh; only the differences between rows are used. A separate calculation
of ionisation potential depression in dense plasma reproduces Stewart--Pyatt in the dilute limit
($3.61$ against $3.57$~eV at $n=10^{21}$~cm$^{-3}$) but becomes unphysical at high density, where a
smeared neutralising background can no longer stand in for explicit neighbouring domains; no
conclusion is drawn from it here. Scripts are available with the source.

\begin{thebibliography}{9}
\bibitem{realqm} C.\ Johnson, \emph{RealQM: Quantum Mechanics as Multiphase 3D Continuum Mechanics},
in~\cite{realqmbook}.
\bibitem{realqmbook} C.\ Johnson, \emph{Real Quantum Mechanics} (Body\&Soul, Vol.\ VI),
\url{https://claes542.github.io/RealMolecule/realquantum3.pdf}.
\bibitem{realnucleus} C.\ Johnson, \emph{RealNucleus: Nuclear Binding as Dual Coulomb Confinement,
without a Strong or Weak Force}, in~\cite{realqmbook}.
\end{thebibliography}

\end{document}
